\begin{pastexam}[label={px:nccu-110-dsa-01}]{政大資科所 110 學年度 第 1 題}
（問答）

(10\%) Prove or disprove the following statements ($n$ is the size of the input array):
\begin{enumerate}[label=\alph*.]
\item (5\%) The time complexity of any sorting algorithm is $O(n^2)$.
\item (5\%) The best-case time complexity of insertion sort is $\theta(n\log n)$.
\end{enumerate}
\end{pastexam}

\begin{pastsoln}
兩個敘述都不成立。 原題 b 的小寫 $\theta$ 指的就是 $\Theta$。

\textbf{a. 不成立}。 要否定「\textbf{任何}排序演算法都是 $O(n^2)$」，只要舉出一個正確的排序演算法，它的時間
不是 $O(n^2)$。

反例：\textbf{排列排序}。 依字典序一個一個產生 $1, 2, \ldots, n$ 的排列 $\pi$，每產生一個就檢查
$A[\pi(1)] \le A[\pi(2)] \le \cdots \le A[\pi(n)]$ 是否成立；成立就照 $\pi$ 輸出，停止。 $n!$ 個排列裡一定有一個
能把 $A$ 排好，所以它一定會停，而且輸出正確。 輸入由大到小排好（鍵互不相同）時，唯一能排好的排列是
$\langle n, n-1, \ldots, 1 \rangle$，它是字典序的最後一個，所以要檢查全部 $n!$ 個排列，每個至少花一步，時間
至少 $n!$。 $n \ge 6$ 時 $(n-1)(n-2) = n^2/2 + (n^2/2 - 3n + 2) \ge n^2/2$，所以
\[
n! \ge n(n-1)(n-2) \ge \frac{n^3}{2}, \qquad \frac{n!}{n^2} \ge \frac{n}{2}.
\]
$n/2$ 會無限變大，不可能有常數 $c$ 使 $n$ 夠大時都有 $n! \le cn^2$。 這個演算法的最壞時間不是 $O(n^2)$，
敘述不成立。

不那麼刻意的例子也有：計數排序（資結書第~15.2.1~節）在鍵的範圍是 $0$ 到 $n^3$ 時要 $\Theta(n + n^3)$。 常見的
比較排序最壞都不超過 $O(n^2)$，但題目說的是「任何」排序演算法。 如果把原句讀成「排序這個\textbf{問題}可以在
$O(n^2)$ 時間內解決」，那是對的（插入排序最壞是 $\Theta(n^2)$）；作答時把這個差別寫出來。

\textbf{b. 不成立}。 插入排序（資結書第~1.3~節）的最好情況是 $\Theta(n)$，不是 $\Theta(n \lg n)$。

\textbf{證明}。 輸入已經由小到大排好時，外層迴圈的每一輪，內層第一次檢查「前一個數是否大於 $\mathit{key}$」
就不成立，不挪動。 每一輪做常數的工作，$n - 1$ 輪共 $\Theta(n)$，所以最好情況的時間 $B(n) \le c_1 n$（$n$ 夠大）。
假設 $B(n) = \Theta(n \lg n)$，就有常數 $c_2 > 0$ 使 $n$ 夠大時 $c_2 n \lg n \le B(n) \le c_1 n$，也就是
$\lg n \le c_1 / c_2$；但 $\lg n$ 會無限變大，矛盾。 所以敘述不成立，正確的說法是插入排序的最好情況是
$\Theta(n)$（任何輸入都至少要把 $n$ 個數各看一次，不會更少）。
\end{pastsoln}
